\chapter{Calf Pain}

\section*{Definition}
Pain localized to the posterior lower leg (gastrocnemius and soleus muscles), ranging from mild soreness to severe debilitating pain.

\section*{Pathophysiology}
Calf pain most commonly arises from musculoskeletal strain (gastrocnemius or soleus tears), but critical vascular causes (DVT, compartment syndrome, arterial occlusion) must be excluded. The calf muscles are particularly vulnerable to strain during activities involving sudden push-off (sprinting, jumping).

\section*{Etiology}
\begin{itemize}
  \item Gastrocnemius or soleus strain (tennis leg)
  \item Deep vein thrombosis (DVT)
  \item Muscle contusion or hematoma
  \item Achilles tendinopathy or rupture
  \item Baker's cyst rupture (pseudothrombophlebitis)
  \item Peripheral arterial disease (claudication)
  \item Compartment syndrome
  \item Cellulitis or soft tissue infection
  \item Electric contusion (charlie horse)
  \item Ruptured popliteal aneurysm
\end{itemize}

\section*{Indications for Evaluation}
Seek care if:
\begin{itemize}
  \item Severe or worsening symptoms
  \item Sudden severe calf pain with swelling, redness, and warmth (DVT concern), or pain after trauma with inability to bear weight
  \item Accompanied by other concerning signs
\end{itemize}

\section*{Related Symptoms}
\begin{itemize}
  \item Leg swelling
  \item Knee pain
  \item Ankle pain
  \item Numbness or tingling
  \item Redness or warmth
\end{itemize}
\textbf{Note:} Always consider serious underlying conditions when evaluating persistent calf pain.

\section*{Self-Care / Home Management}
\begin{itemize}
  \item Rest and avoid activities that may aggravate the condition.
  \item Maintain adequate hydration and a balanced diet.
  \item Over-the-counter remedies may provide symptomatic relief (consult a pharmacist or doctor).
  \item Monitor symptoms and seek professional advice if they persist or worsen.
  \item Avoid self-medicating with prescription medications without medical guidance.
\end{itemize}

\section*{Diagnostic Approach}
\begin{itemize}
  \item A thorough history and physical examination are the first steps in evaluation.
  \item Laboratory tests (blood work, urinalysis) may be ordered based on clinical suspicion.
  \item Imaging studies (X-ray, ultrasound, CT, MRI) may be indicated when structural pathology is suspected.
  \item Specialist referral is arranged when the diagnosis remains uncertain or requires specific expertise.
\end{itemize}

\section*{Treatment / Management Overview}
\begin{itemize}
  \item Treatment targets the underlying cause identified through diagnostic evaluation.
  \item Supportive care and symptom management are provided while awaiting definitive therapy.
  \item Medications, lifestyle modifications, and physical therapies may be prescribed as appropriate.
  \item Follow-up ensures treatment effectiveness and allows timely adjustments to the care plan.
\end{itemize}

\clearpage